\section{Introdução}\label{sec:introducao}

Ferramentas de PLN têm sido aplicadas à revisão de textos acadêmicos há pelo
menos uma década \cite{lecun2015}. Neste trabalho, o processamento de linguagem
natural (PLN) é usado para apoiar autores que escrevem em português.

A revisão manual de manuscritos é custosa devido à fatores como o volume
crescente de submissões. À partir de 2020, diversos periódicos passaram a
exigir verificações automáticas de similaridade \cite{devlin2019}.

A idéia central deste trabalho é avaliar o texto a nível de parágrafo e de
seção, apontando problemas que um revisor humano costuma encontrar.

O objetivo deste trabalho é descrever uma ferramenta de auditoria que lê os
arquivos-fonte de um artigo, identifica problemas de estrutura, de referências
e de escrita, e apresenta cada achado com o arquivo, a linha e uma sugestão de
correção, sem jamais alterar os arquivos originais do autor.

Trabalhos anteriores \cite{silva2099} propuseram verificadores para a língua
inglesa, mas poucos tratam das normas brasileiras.
% TODO revisar esta seção antes de enviar
